% Generated by trazabilidad/matriz.py
\multicolumn{4}{@{}l}{\textit{CORE y Cuentas}} \\
\nm{fn\_asiento\_cuadra} (trigger de restricción diferido) & Un asiento tiene al menos una línea en el debe y otra en el haber, y sus sumas coinciden; se verifica al confirmar la transacción. \textit{Pruebas: M20a, M20b, F24.} & CORE y Cuentas (RN-COR-002) \parencites{BNA2026Core}[p.~16]{BNA2026Cuentas} & Nuevo \\
\nm{fn\_solo\_insercion} (triggers por fila y por sentencia) & Asientos, historiales, intentos, auditoría y aplicaciones de pago no admiten \texttt{UPDATE}, \texttt{DELETE} ni \texttt{TRUNCATE}; las autorizaciones no se borran. \textit{Pruebas: F11a, F11b, F11c, F11d, F11e, F11f, F13, F20.} & CORE, Tarjetas (RN-SEG-002, RN-ATM-003) y Desembolso; \texttt{TRUNCATE}, aplicaciones de pago e historial de la solicitud, decisión de este informe \parencites{BNA2026Core}[p.~23]{BNA2026Tarjetas}[p.~13]{BNA2026Desembolso} & Modelo \\
\nm{fn\_fecha\_negocio} & El día de negocio se corta a las 00:00 de Arequipa. \textit{Pruebas: F12l, F12m, F12n.} & Tarjetas, RN-TOP-001, y Desembolso \parencites[p.~15]{BNA2026Tarjetas}[p.~4]{BNA2026Desembolso} & Nuevo \\
\nm{fn\_movimiento\_validar}, \nm{fn\_movimiento\_aplicar}, \nm{fn\_consumo\_sumar} & La cuenta admite débitos o créditos según su estado; el asiento está en la moneda de la cuenta; un débito sin autorización no supera el saldo ni el tope del canal. Bloquea la cuenta, actualiza los saldos y el consumo del día, y vincula el movimiento a su autorización; el extorno devuelve el consumo al día en que se comprometió. \textit{Pruebas: M04, M10a, M10b, M22, M24a, F21a, F21b, F22, F12m, C01, C02.} & Cuentas, CORE y Tarjetas (RN-TOP-001, RN-BD-002, C07) \parencites[pp.~16, 17, 25 y~32]{BNA2026Cuentas}{BNA2026Core}[pp.~15 y~28]{BNA2026Tarjetas} & Modelo \\
\nm{fn\_movimiento\_inmutable} & Un movimiento solo cambia de APROBADO a EXTORNADO. \textit{Pruebas: F23.} & CORE \parencite{BNA2026Core} & Nuevo \\
\addlinespace
\multicolumn{4}{@{}l}{\textit{Tarjetas}} \\
\nm{fn\_tarjeta\_emitir} & Solo se emite tarjeta a un cliente ACTIVO (RF-TAR-005). \textit{Pruebas: F01.} & Tarjetas, \nm{trg\_tarjeta\_bi} \parencite[pp.~23 y~31]{BNA2026Tarjetas} & Portado \\
\nm{fn\_tarjeta\_emitida} & La emisión queda en el historial y en la bandeja hacia Seguridad (RN-SEG-002). \textit{Pruebas: F02.} & Tarjetas, \nm{trg\_tarjeta\_ai} \parencite[pp.~22 y~23]{BNA2026Tarjetas} & Portado \\
\nm{fn\_historial\_validar}, \nm{fn\_historial\_aplicar}, \nm{fn\_tarjeta\_estado\_via\_historial} & La cancelada no cambia y la vencida solo se cancela (RN-TAR-003); solo un operador o Seguridad desbloquean, y al desbloquear los intentos vuelven a cero (RN-TAR-004). El estado de la tarjeta sigue a su historial y cada cambio se avisa a Seguridad. \textit{Pruebas: F03, F04, F05, F06, F07.} & Tarjetas, \nm{trg\_tarjeta\_bu} y \nm{trg\_tarjeta\_au}; el cambio solo por el historial, decisión de este informe \parencite[pp.~23 y~29]{BNA2026Tarjetas} & Portado \\
\nm{fn\_intento\_contar} & El tercer intento fallido de la credencial o de la matriz bloquea la tarjeta; un intento exitoso reinicia el contador (RN-TAR-004). \textit{Pruebas: F08, F09, F10.} & Tarjetas, \nm{trg\_intento\_ai} \parencite[pp.~19 y~23]{BNA2026Tarjetas} & Portado \\
\nm{fn\_autorizacion\_validar}, \nm{fn\_validar\_reglas\_autorizacion}, \nm{fn\_autorizacion\_aplicar} & Para autorizar, la tarjeta está operativa y vigente, la terminal y el canal están activos, el cajero no opera cuentas mancomunadas conjuntas (RN-ATM-005), el canal opera en la moneda (RN-BD-002) y el débito no supera el tope (RN-TOP-001) ni el saldo, con la cuenta bloqueada. Al autorizar, compromete el tope y retiene el saldo; al extornar sin movimiento, los libera. Controla la secuencia de estados (RN-ATM-007). \textit{Pruebas: M23, M24b, M24c, F12a, F12b, F12c, F12d, F12e, F12f, F12g, F12h, F12i, F12j, F12k, F12l, F12n, C03.} & Tarjetas, \nm{trg\_autorizacion\_bi/bu} y \nm{sp\_validar\_reglas\_autorizacion} \parencite[pp.~15, 21, 23, 26 y~27]{BNA2026Tarjetas} & Portado \\
\nm{fn\_cliente\_estado\_tarjetas} & Si Clientes bloquea o suspende al cliente, se bloquean sus tarjetas; si lo da de baja, se cancelan (INT-T07). \textit{Pruebas: F14a, F14b.} & Tarjetas, \nm{trg\_cliente\_ref\_au}, y Clientes \parencites[pp.~23 y~31]{BNA2026Tarjetas}[p.~8]{BNA2026Clientes} & Portado \\
\nm{bloquear\_tarjeta()} & Bloquea la tarjeta registrando su motivo y su origen en el historial. \textit{Pruebas: M26.} & Tarjetas, RN-SEG-002 \parencite[pp.~18 y~22]{BNA2026Tarjetas} & Portado \\
\nm{marcar\_tarjetas\_vencidas()} & Marca como VENCIDA cada tarjeta cuya fecha ya pasó (RN-TAR-003). \textit{Pruebas: F15.} & Tarjetas, \nm{ev\_marcar\_tarjetas\_vencidas} \parencite[pp.~23 y~32]{BNA2026Tarjetas} & Portado \\
\nm{v\_tarjeta\_resumen}, \nm{v\_consumo\_diario}, \nm{v\_tope\_restante\_hoy}, \nm{v\_credenciales\_en\_riesgo} & Número enmascarado y vencimiento MM/AA (PCI-DSS); consumo por tarjeta, canal, moneda y día; tope disponible por tarjeta y canal (RN-TOP-001); intentos acumulados y restantes (RN-TAR-004). \textit{Pruebas: M25, F25, F26, F27.} & Tarjetas, objetos programados y consultas C05 y C08 \parencite[pp.~23, 24, 27 y~28]{BNA2026Tarjetas} & Portado \\
\addlinespace
\multicolumn{4}{@{}l}{\textit{Solicitud, Evaluación, Desembolso y Cobranza}} \\
\nm{fn\_orden\_validar} & La cuenta destino pertenece al cliente de la solicitud y tiene su moneda; el importe no supera el límite aprobado (decisión de este informe); hay una sola orden CONFIRMADA por solicitud, en todas sus rondas. \textit{Pruebas: M09a, M09b, M09c, M14b.} & Desembolso y Evaluación \parencites[pp.~11 y~12]{BNA2026Desembolso}[p.~8]{BNA2026Evaluacion} & Modelo \\
\nm{v\_ordenes\_inciertas}, \nm{v\_entregas\_pendientes}, \nm{v\_conciliacion\_diaria} & Órdenes con resultado incierto en el CORE; entregas a Cobranza que superan dos reintentos; importes ordenados frente a abonados, por fecha y moneda. \textit{Pruebas: F28, F29, F30.} & Desembolso, vistas de apoyo operativo; reintentos, documento de tecnologías \parencites[p.~12]{BNA2026Desembolso}[línea~92]{BNA2026Stack} & Portado \\
\nm{fn\_credito\_validar} & El crédito nace de una orden CONFIRMADA (decisión I15) y su calendario pertenece a la solicitud de la orden y es OFICIAL (decisión de este informe). \textit{Pruebas: M14a, F16a, F16b, F16c.} & Cobranza (RF-COB-015) y Solicitud \parencites[p.~2]{BNA2026Cobranza}[p.~4]{BNA2026Solicitud} & Modelo \\
\nm{fn\_aplicacion\_validar}, \nm{fn\_aplicacion\_aplicar} & Solo se aplica un pago CONFIRMADO, a cuotas de su crédito, sin superar el saldo de la cuota ni el monto del pago. Actualiza el saldo y el estado de la cuota y cancela el crédito cuando todas sus cuotas quedan pagadas. \textit{Pruebas: M05, M13a, M13b, F17a, F17b, F19.} & Cobranza \parencite[pp.~3 y~8]{BNA2026Cobranza} & Modelo \\
\nm{aplicar\_pago()} & Reparte el saldo libre de un pago de la cuota más antigua a la siguiente. \textit{Pruebas: F18.} & Cobranza \parencite[p.~3]{BNA2026Cobranza} & Portado \\
\addlinespace
\multicolumn{4}{@{}l}{\textit{Índices}} \\
\stat{consultas} índices de consulta & Movimientos, autorizaciones, intentos e historial por fecha; órdenes por estado y fecha; eventos por enviar; auditoría por \nm{correlation\_id}. \textit{Pruebas: F31.} & Tarjetas y Desembolso \parencites[p.~30]{BNA2026Tarjetas}[p.~12]{BNA2026Desembolso} & Portado \\
